\chapter{Background and Related Work}\label{ch:background}
The relevant literature concerns both the physical meaning of an image-derived corrosion measure and the evaluation of learned relationships. This chapter uses a focused selection of verified sources to establish those distinctions; it is not a systematic literature review.

\section{Predecessor experimental and imaging work}
Hossain's thesis describes the experimental ferrocement specimens, exposure procedures, photographic record, mechanical testing, and microscopic investigation \citep{hossain_thesis}. In that work, four-point bending is performed after ageing; the microscopic assessment examines longitudinal wires in the layer near the cracked or exposed surface. These measurements provide the physical context for the terminal endpoints used here. A selected cross-sectional area-loss measurement does not constitute an observed internal-damage trajectory over the exposure period.

The related conference paper develops an image-analysis and classification workflow for ferrocement surface corrosion \citep{artiste2025}. Its RGB-based analysis and three-category classification belong to that external predecessor study. The present repository contains different computational generations and a later four-class preparation branch. Their class schemes, data handling, and trained-model status must remain separate. The predecessor's classifier results cannot be used as results of the current four-class branch.

The contribution of the present activity is therefore a secondary modelling and evaluation study on the inherited observations. Reusing this experimental basis allows representation and evaluation comparisons, but does not provide an independent replication of the physical experiment.

\section{Image-based corrosion and structural monitoring}
Atha and Jahanshahi evaluate convolutional-network approaches to corrosion detection \citep{atha2018}. Their work motivates visual feature learning for detecting corrosion evidence. It does not establish that a detector's score estimates internal metal loss or terminal capacity in ferrocement. Those are different targets requiring their own ground truth and validation.

Spencer, Hoskere, and Narazaki distinguish inspection applications from monitoring applications in their review of computer vision for civil infrastructure \citep{spencer2019}. That distinction is useful here: identifying visible damage and measuring structural response need different observation models. Repeated photographs can support an inspection time series, but the availability of a time series alone does not make its downstream structural interpretation longitudinally validated.

Frozen residual-network embeddings provide one intermediate representation between colour rules and fully task-specific deep learning. Residual networks were developed to facilitate the optimization of deep visual models \citep{he2016}; in this activity, an existing network supplies fixed image features and a separately trained regression head predicts the visible label. The saved experiment is evidence about that particular transfer-learning arrangement, not about the universal superiority of deep or handcrafted representations.

\section{Evaluation must match the prediction question}
Repeated observations create a dependence structure. Roberts et al. discuss cross-validation for structured data and show why the blocking strategy must reflect the intended prediction problem \citep{roberts2017}. Their methodological discussion also matters for interpretation: withholding a group can reduce overlap in the predictor space. Here, specimen exclusion prevents a specimen's image history from appearing in both fitting and evaluation, while campaign exclusion additionally removes a whole design/exposure combination.

Kapoor and Narayanan connect leakage risks to the validity of scientific claims made from machine-learning performance \citep{kapoor2023}. In this study the essential information-flow questions are concrete: whether a specimen crosses partitions, whether terminal targets enter the feature matrix, whether preprocessing uses held-out data, and whether model-generated structural estimates are subsequently mistaken for observations. A clean specimen split addresses one of these questions, not all of them.

Model selection creates an additional boundary. Cawley and Talbot show that optimizing a finite-sample evaluation criterion can introduce selection bias \citep{cawley2010}. Accordingly, a selected winner's reported fold score should not be presented as an unbiased estimate of the full model-selection procedure. This is especially pertinent when image-family differences are small and the number of independent structural observations is limited.

\section{Position of this study}
The evidence reviewed above motivates a specific contribution: testing whether visible-corrosion representations support incremental structural prediction under realistic information and grouping constraints. The study supplies a bounded empirical case, not a new general-purpose model. Its value lies in explaining how apparently encouraging pooled scores change when the target, observation unit, and evaluation population are specified. The next chapter makes that observation structure explicit.
